\begin{proposition}[The Chow obstruction on an explicit fixed open]
\label{prop:B-fixed-open-ramification-obstruction}
\status{Proved exact comparison formula; no colimit vanishing or survival assertion}
Let \(k=\F_q\), let \(\ell\ne\operatorname{char}k\), and put
\(Y=\mathbb P^4_k\).  Let \(D_1,D_2\subset Y\) be smooth
hypersurfaces of positive degree meeting transversely in a smooth
surface \(S\).  Write \(U=Y\setminus(D_1\cup D_2)\) and
\(\Gamma=\Gal(\bar k/k)\).  Assume the divisor Tate conjecture
for \(S\).  Define
\[
\begin{aligned}
 V&=\operatorname{CH}^1(S)_\Q\otimes\Q_\ell,\\
 C&=\bigl(\operatorname{CH}^2(D_1)_\Q
             \oplus\operatorname{CH}^2(D_2)_\Q\bigr)
                 \otimes\Q_\ell,\\
 G:V&\longrightarrow C,
       &G(z)&=((i_1)_*z,-(i_2)_*z),
\end{aligned}
\]
where \(i_a:S\hookrightarrow D_a\) are the inclusions.  Let
\[
 \operatorname{cl}_C:C\longrightarrow
 H^4(\bar D_1,\Q_\ell(2))\oplus H^4(\bar D_2,\Q_\ell(2))
\]
be the ordinary cycle map.  There is an isomorphism
\[
\begin{split}
 \operatorname{coker}\bigl[
   H_M^4(U,\Q(3))\otimes\Q_\ell
       &\longrightarrow H^4(\bar U,\Q_\ell(3))^\Gamma
                         \bigr]\\
       &\simeq\operatorname{im}G\cap\ker\operatorname{cl}_C.
\end{split}
\]
Here \(H_M\) is ordinary rational motivic cohomology and the
right-hand cohomology is continuous \(\ell\)-adic cohomology.
The identification is canonical after fixing the displayed
Mayer--Vietoris orientation and the usual cycle maps.
\end{proposition}
\begin{proof}
Put \(D=D_1\cup D_2\).  The Mayer--Vietoris triangle for
supports and purity give
\[
\begin{split}
 R\Gamma(\bar S,\Q_\ell(1))[-4]
   &\longrightarrow
      R\Gamma(\bar D_1,\Q_\ell(2))[-2]
        \oplus R\Gamma(\bar D_2,\Q_\ell(2))[-2]\\
   &\longrightarrow R\Gamma_{\bar D}(\bar Y,\Q_\ell(3)).
\end{split}
\]
For the support triangle, apply ordinary open Mayer--Vietoris
to the cover
\((Y\setminus D_1)\cup(Y\setminus D_2)=Y\setminus S\),
compare with restriction from \(Y\), and take the three
mapping cones in functorial cochain models.  The same construction
works motivically; no descent across the singular closed immersion
\(D\hookrightarrow Y\) is assumed.  The purity shifts are those
for the regular immersions of codimensions two and one,
respectively~\cite[Theorem~3.1.1, p.~323]{c_riou2014}.

The relevant support sequence is
\[
\begin{aligned}
 H^1(\bar S,\Q_\ell(1))
 &\longrightarrow\bigoplus_{a=1}^2H^3(\bar D_a,\Q_\ell(2))
 \longrightarrow H^5_{\bar D}(\bar Y,\Q_\ell(3))\\
 &\longrightarrow H^2(\bar S,\Q_\ell(1))
 \longrightarrow\bigoplus_{a=1}^2H^4(\bar D_a,\Q_\ell(2)).
\end{aligned}
\]
The first two terms are pure of weight \(-1\).  Hence the
kernel of the map from the support group to the last Gysin
kernel has \(F-1\) invertible, where \(F\) is geometric
Frobenius.  Both its \(\Gamma\)-invariants and its first
\(\Gamma\)-cohomology vanish.  Taking invariants therefore
identifies the support invariants with the Gysin kernel on
invariant subspaces.  This uses the proper smooth weight
theorem, without a separate semisimplicity assumption
\cite[Corollary~3.3.9, p.~207]{c_deligneweil1980}.

Next, localization gives
\[
 H^4(\bar Y,\Q_\ell(3))\longrightarrow
 H^4(\bar U,\Q_\ell(3))\longrightarrow
 H^5_{\bar D}(\bar Y,\Q_\ell(3))\longrightarrow
 H^5(\bar Y,\Q_\ell(3))=0.
\]
The image of the first term is pure of weight \(-2\), so
the same invariant and coinvariant argument identifies
\(H^4(\bar U,\Q_\ell(3))^\Gamma\) with the support
invariants.  Divisor Tate on \(S\), together with rational
Kummer injectivity for its Picard group, identifies
\[
 V\simeq H^2(\bar S,\Q_\ell(1))^\Gamma.
\]
Thus the continuous target in the statement is exactly
\(\ker(\operatorname{cl}_C G)\subset V\).

For ordinary motivic cohomology, the analogous support sequence
has the segment
\[
 H^5_D(Y,\Q(3))\longrightarrow
 H_M^2(S,\Q(1))\longrightarrow
 H_M^4(D_1,\Q(2))\oplus H_M^4(D_2,\Q(2)).
\]
The last three groups are the corresponding rational Chow
groups.  Its image in \(\operatorname{CH}^1(S)_\Q\)
is consequently the kernel of the rational signed Gysin map.
Motivic localization also gives
\[
 H_M^4(U,\Q(3))\longrightarrow H^5_D(Y,\Q(3))
                   \longrightarrow H_M^5(\mathbb P^4_k,\Q(3)).
\]
The last group is zero: the projective-bundle formula expresses it as
\[
 \bigoplus_{r=0}^4 H_M^{5-2r}(k,\Q(3-r))=0,
\]
by the ordinary finite-field motivic calculation.  These are
the usual motivic purity, localization and projective-bundle
maps~\cite[(14.5.1), Theorem~15.12, Corollary~15.13,
Theorem~15.15]{BMVW2006}.
The ordinary cycle map commutes with these maps.  Since
\(\Q_\ell\) is flat over \(\Q\), its image in \(V\)
is therefore exactly \(\ker G\).  Finally, applying \(G\)
induces the elementary isomorphism
\[
 \ker(\operatorname{cl}_C G)/\ker G
       \simeq\operatorname{im}G\cap\ker\operatorname{cl}_C,
\]
which proves the assertion.
\end{proof}

The formula locates a concrete obstruction: divisor Tate on the
surface supplies an algebraic representative, but the outgoing
Gysin relation is initially only a cohomological relation on the
threefold boundary components.  It need not yet be a rational
Chow relation.  We have neither proved that this intersection
vanishes nor constructed a nonzero element of it.
Moreover, failure of motivic surjectivity on this fixed \(U\)
does not show survival under every further shrinking: a class
could become zero on a smaller open without having had a motivic
lift on the original one.  The proposition thus neither solves
nor refutes the projective-space colimit gate.  It assumes only
the displayed smooth transverse pair, without asserting that
every open admits such a compactification.
